% =====================================================================
%  CARNET D'ENTRAÎNEMENT — FICHE 3 : Trigonométrie
%  Thème 2 — Angles associés
%  TUTORIEL
% =====================================================================
\documentclass[11pt,a4paper]{article}
\usepackage[utf8]{inputenc}
\usepackage[T1]{fontenc}
\usepackage{lmodern}
\usepackage[french]{babel}
\usepackage{amsmath,amssymb,mathtools}
\usepackage{icomma}
\usepackage{xcolor}
\usepackage{colortbl}
\usepackage{tikz}
\usepackage[most]{tcolorbox}
\usepackage{enumitem}
\usepackage{array}
\usepackage{booktabs}
\usepackage[a4paper,margin=2.2cm]{geometry}
\usetikzlibrary{arrows.meta,calc,angles,quotes,decorations.markings}

\definecolor{primary}{RGB}{214,51,108}
\definecolor{primarydark}{RGB}{140,20,64}
\definecolor{defcol}{RGB}{0,114,178}
\definecolor{excol}{RGB}{0,135,90}
\definecolor{methcol}{RGB}{90,90,100}
\definecolor{attncol}{RGB}{200,40,55}
\definecolor{sincol}{RGB}{205,40,55}
\definecolor{coscol}{RGB}{20,110,200}

\DeclareMathOperator{\tg}{tg}
\DeclareMathOperator{\cotg}{cotg}
\newcommand{\degr}{^{\circ}}
\newcommand{\Z}{\mathbb{Z}}

\tcbset{
  boxbase/.style={enhanced,breakable,boxrule=0.9pt,arc=2.5pt,
     left=3mm,right=3mm,top=2.3mm,bottom=2.3mm,
     fonttitle=\bfseries,coltitle=white,toptitle=0.6mm,bottomtitle=0.6mm},
}
\newtcolorbox{idee}[1][]{boxbase,colback=defcol!6,colframe=defcol,
  title={\textsc{L'essentiel}\ifx\relax#1\relax\relax\else\textmd{ : \itshape #1}\fi}}
\newtcolorbox{exemple}[1][]{boxbase,colback=excol!6,colframe=excol,
  title={\textsc{Exemple}\ifx\relax#1\relax\relax\else\textmd{ : \itshape #1}\fi}}
\newtcolorbox{methode}[1][]{boxbase,colback=methcol!8,colframe=methcol,sharp corners,
  title={\textsc{Méthode}\ifx\relax#1\relax\relax\else\textmd{ --- \itshape #1}\fi}}
\newtcolorbox{piege}[1][]{boxbase,colback=attncol!7,colframe=attncol,
  title={\textsc{Piège à éviter}\ifx\relax#1\relax\relax\else\textmd{ : \itshape #1}\fi}}

\newcommand{\heading}[1]{\smallskip\noindent{\large\bfseries\color{primarydark}#1}\par\smallskip}

\pagestyle{empty}
\setlength{\parindent}{0pt}
\setlength{\parskip}{2pt}

\begin{document}

\begin{center}
{\LARGE\bfseries\color{primary}Thème 2 — Angles associés}\\[2pt]
{\color{primarydark}\textsc{Carnet d'entraînement — Fiche 3 : Trigonométrie \quad•\quad Tutoriel}}
\end{center}
\hrule height 1pt
\medskip

Deux angles sont \textbf{associés} lorsque leurs points sur le cercle sont \textbf{symétriques} (par un
axe ou par le centre) : leurs nombres trigonométriques sont alors les mêmes \emph{au signe près}, ou
\emph{échangés} ($\sin\leftrightarrow\cos$). Objectif : \textbf{ramener n'importe quel angle à un angle
connu} du premier quadrant.

\heading{1. Les cinq familles à connaître}
Dans tout ce qui suit, $\alpha$ est un angle \emph{quelconque}.

\begin{center}
\renewcommand{\arraystretch}{1.3}
\begin{tabular}{|>{\columncolor{primary!10}}l|c|c|c|c|}
\hline
\rowcolor{primary!20}
\textbf{Famille (symétrie)} & \textbf{angle} & $\cos$ & $\sin$ & $\tg$\\
\hline
Opposés \emph{(axe des cos)} & $-\alpha$ & $\cos\alpha$ & $-\sin\alpha$ & $-\tg\alpha$\\
\hline
Supplémentaires \emph{(axe des sin)} & $\pi-\alpha$ & $-\cos\alpha$ & $\sin\alpha$ & $-\tg\alpha$\\
\hline
Anti-suppl. \emph{(centre $O$)} & $\pi+\alpha$ & $-\cos\alpha$ & $-\sin\alpha$ & $\tg\alpha$\\
\hline
Complémentaires \emph{(droite $y=x$)} & $\frac{\pi}{2}-\alpha$ & $\sin\alpha$ & $\cos\alpha$ & $\cotg\alpha$\\
\hline
Anti-complém. & $\frac{\pi}{2}+\alpha$ & $-\sin\alpha$ & $\cos\alpha$ & $-\cotg\alpha$\\
\hline
\end{tabular}
\end{center}

\begin{idee}[deux réflexes qui résument tout]
$\bullet$ \textbf{Autour de $\pi$} (c'est-à-dire $-\alpha,\ \pi-\alpha,\ \pi+\alpha$) : la
\emph{fonction ne change pas} ($\cos$ reste $\cos$, $\sin$ reste $\sin$), seul le \textbf{signe} peut
changer.\\
$\bullet$ \textbf{Autour de $\frac{\pi}{2}$} (c'est-à-dire $\frac{\pi}{2}\pm\alpha$) : la fonction
\emph{s'échange} ($\sin\leftrightarrow\cos$, $\tg\leftrightarrow\cotg$), et le signe peut changer.\\
Pour le signe : place $\alpha$ « petit et positif » et regarde le quadrant de l'angle associé.
\end{idee}

\heading{2. Lire une symétrie sur le cercle}
Ci-dessous : les \textbf{opposés} $\alpha$ et $-\alpha$ (miroir par l'axe horizontal : même $\cos$,
$\sin$ opposé) et les \textbf{supplémentaires} $\alpha$ et $\pi-\alpha$ (miroir par l'axe vertical :
même $\sin$, $\cos$ opposé).

\begin{center}
\begin{tikzpicture}[scale=1.35,line cap=round,>=Stealth]
  \draw[primary,line width=1pt] (0,0) circle (1);
  \draw[->,black!70] (-1.35,0)--(1.4,0) node[right]{\scriptsize$\cos$};
  \draw[->,black!70] (0,-1.35)--(0,1.4) node[above]{\scriptsize$\sin$};
  % alpha au quadrant I
  \draw[primarydark,line width=1pt] (0,0)--(50:1);
  \fill[primarydark] (50:1) circle (0.028) node[above right]{\scriptsize$\alpha$};
  % supplementaire pi - alpha (quadrant II)
  \draw[coscol,line width=1pt] (0,0)--(130:1);
  \fill[coscol] (130:1) circle (0.028) node[above left]{\scriptsize$\pi-\alpha$};
  % oppose -alpha (quadrant IV)
  \draw[excol,line width=1pt] (0,0)--(-50:1);
  \fill[excol] (-50:1) circle (0.028) node[below right]{\scriptsize$-\alpha$};
  % projections
  \draw[black!30,dashed] (50:1)--(-50:1);
  \draw[black!30,dashed] (50:1)--(130:1);
  \fill (0,0) circle (0.02);
\end{tikzpicture}
\end{center}

\begin{exemple}[même sinus, même tangente]
Deux supplémentaires ont le même sinus : $\sin150\degr=\sin(180\degr-30\degr)=\sin30\degr=\dfrac12$.
Deux anti-supplémentaires ont la même tangente :
$\tg210\degr=\tg(180\degr+30\degr)=\tg30\degr=\dfrac{\sqrt3}{3}$ (période $\pi$).
\end{exemple}

\heading{3. Ramener n'importe quel angle au premier quadrant}
\begin{methode}[réduction d'un angle quelconque]
\begin{enumerate}[leftmargin=1.6em,topsep=1pt,itemsep=1pt]
  \item \textbf{Un tour} : ajoute/retranche $2\pi$ (ou $360\degr$) pour tomber dans $[0;2\pi[$.
  \item \textbf{Quadrant $\Rightarrow$ signe} : repère le quadrant, il donne le signe du résultat.
  \item \textbf{Angle de référence} : écris l'angle sous la forme $\pi\pm\beta$, \ $-\beta$ \ ou
        \ $\frac{\pi}{2}\pm\beta$ avec $\beta$ \emph{aigu} (le référent du 1\textsuperscript{er} quadrant).
  \item \textbf{Applique la famille} correspondante, puis recolle le signe de l'étape 2.
\end{enumerate}
\end{methode}

\begin{exemple}[valeur numérique et expression littérale]
$\dfrac{2\pi}{3}=\pi-\dfrac{\pi}{3}$ \emph{(supplémentaires)} $\Rightarrow
\cos\dfrac{2\pi}{3}=-\cos\dfrac{\pi}{3}=-\dfrac12$ (cohérent : $120\degr$ est au quadrant II, où
$\cos<0$).\quad
De même $\cos\!\left(\dfrac{\pi}{2}+\alpha\right)=-\sin\alpha$ \emph{(anti-complém.)} et
$\dfrac{\sin(\pi-\alpha)}{\cos(-\alpha)}=\dfrac{\sin\alpha}{\cos\alpha}=\tg\alpha$.
\end{exemple}

\begin{piege}[les trois erreurs qui coûtent des points]
$\bullet$ Confondre \emph{supplémentaires} ($\pi-\alpha$, même $\sin$) et \emph{anti-supplémentaires}
($\pi+\alpha$, même $\tg$).\quad
$\bullet$ Oublier que $\frac{\pi}{2}\pm\alpha$ \textbf{échange} $\sin$ et $\cos$ (ne pas laisser $\cos$
donner un $\cos$).\quad
$\bullet$ Se tromper de signe : dans le doute, prends $\alpha=30\degr$ et vérifie numériquement.
\end{piege}

\end{document}
